\section{Projective Covers over Artinian Rings}

\begin{lemma}\label{lem:primitive-decomposition-semisimple-artinian}
    Let $A$ be a semisimple Artinian ring. Then the following statements hold.

    \begin{enumerate}
        \item There exists a primitive idempotent decomposition
              \[
                  1=e_1+\cdots+e_n .
              \]
              This decomposition corresponds to a decomposition of the left regular module
              \[
                  {}_A A = Ae_1\oplus \cdots \oplus Ae_n,
              \]
              where each $Ae_i$ is a simple left $A$-module.

        \item For every simple left $A$-module $S$, there exists a primitive idempotent
              $e\in A$ such that
              \[
                  S\cong Ae .
              \]

        \item For every simple left $A$-module $S$, there exists some $1\le i\le n$
              such that
              \[
                  e_iS\neq 0.
              \]
              Moreover, if $T$ is a simple left $A$-module with $T\not\cong S$, then
              \[
                  e_iT=0.
              \]
    \end{enumerate}
\end{lemma}

\begin{proof}
    \textbf{(1)}
    Since $A$ is semisimple Artinian, the left regular module ${}_A A$ is semisimple and
    Artinian. Hence it is a finite direct sum of simple left ideals:
    \[
        {}_A A = I_1\oplus \cdots \oplus I_n .
    \]
    By the correspondence between direct sum decompositions of ${}_A A$ and idempotent
    decompositions of $1$, there exist pairwise orthogonal idempotents
    \[
        e_1,\dots,e_n\in A
    \]
    such that
    \[
        1=e_1+\cdots+e_n
    \]
    and
    \[
        I_i=Ae_i
        \qquad (1\le i\le n).
    \]
    Since each $I_i$ is simple, each $Ae_i$ is simple and hence indecomposable. Therefore
    each $e_i$ is primitive.

    \medskip

    \textbf{(2)}
    Let $S$ be a simple left $A$-module. Choose $0\neq s\in S$. Then
    \[
        S=As,
    \]
    so the map
    \[
        A\longrightarrow S,
        \qquad
        a\longmapsto as
    \]
    is a surjective $A$-module homomorphism.

    Using the decomposition from part \textup{(1)}, we have
    \[
        A=Ae_1\oplus\cdots\oplus Ae_n.
    \]
    Since the map $A\to S$ is nonzero and surjective, its restriction to at least one
    summand $Ae_i$ is nonzero. Thus there exists $i$ such that we have a nonzero
    homomorphism
    \[
        Ae_i\longrightarrow S.
    \]
    Because both $Ae_i$ and $S$ are simple, this nonzero homomorphism is an isomorphism.
    Therefore
    \[
        S\cong Ae_i.
    \]
    Taking $e=e_i$, we obtain the desired primitive idempotent.

    \medskip

    \textbf{(3)}
    Let $S$ be a simple left $A$-module. Since
    \[
        1=e_1+\cdots+e_n,
    \]
    we have
    \[
        S=1S=(e_1+\cdots+e_n)S=e_1S+\cdots+e_nS.
    \]
    Since $S\neq 0$, there exists some $i$ such that
    \[
        e_iS\neq 0.
    \]

    For this $i$, we claim that if $T$ is a simple left $A$-module with $T\not\cong S$,
    then
    \[
        e_iT=0.
    \]
    Indeed, by the standard isomorphism
    \[
        \operatorname{Hom}_A(Ae_i,T)\cong e_iT,
        \qquad
        \varphi\longmapsto \varphi(e_i),
    \]
    it is enough to show that
    \[
        \operatorname{Hom}_A(Ae_i,T)=0.
    \]
    Since $e_iS\neq 0$, the same isomorphism gives
    \[
        \operatorname{Hom}_A(Ae_i,S)\cong e_iS\neq 0.
    \]
    Thus there is a nonzero homomorphism
    \[
        Ae_i\longrightarrow S.
    \]
    Because $Ae_i$ and $S$ are simple, this homomorphism is an isomorphism. Hence
    \[
        Ae_i\cong S.
    \]
    If $T\not\cong S$, then $T\not\cong Ae_i$. Therefore, by Schur's lemma,
    \[
        \operatorname{Hom}_A(Ae_i,T)=0.
    \]
    Consequently,
    \[
        e_iT=0.
    \]
\end{proof}

\begin{theorem}\label{thm:projective-cover-simple}
    Let $A$ be an Artinian ring, and let $S$ be a simple left $A$-module. Then the following hold.

    \begin{enumerate}
        \item There exists an indecomposable projective left $A$-module $P_S$ such that
              \[
                  P_S/\operatorname{Rad}(P_S)\cong S.
              \]
              Moreover, $P_S$ is of the form
              \[
                  P_S=Af,
              \]
              where $f$ is a primitive idempotent of $A$.

        \item The idempotent $f$ has the following property:
              \[
                  fS\neq 0,
              \]
              and if $T$ is any simple left $A$-module with $T\not\cong S$, then
              \[
                  fT=0.
              \]

        \item The module $P_S$ is the projective cover of $S$.
    \end{enumerate}
\end{theorem}

\begin{proof}
    Let
    \[
        \overline{A}:=A/J(A).
    \]
    Since $A$ is Artinian, the ideal $J(A)$ is nilpotent, and $\overline{A}$ is a semisimple
    Artinian ring.

    By the structure theory of semisimple Artinian rings, there exists a primitive idempotent
    \[
        e\in \overline{A}
    \]
    such that
    \[
        eS\neq 0.
    \]
    Equivalently, the simple $\overline{A}$-module $\overline{A}e$ is isomorphic to $S$.

    Since $J(A)$ is nilpotent, idempotents lift modulo $J(A)$. Hence there exists an
    idempotent $f\in A$ such that
    \[
        f+J(A)=e.
    \]
    Moreover, because $e$ is primitive, the lift $f$ is also primitive.

    Now define
    \[
        P_S:=Af.
    \]
    Since $f$ is an idempotent, $Af$ is a direct summand of the left regular module
    ${}_A A$. Indeed,
    \[
        A=Af\oplus A(1-f).
    \]
    Hence $Af$ is projective. Since $f$ is primitive, $Af$ is indecomposable. Therefore
    $P_S=Af$ is an indecomposable projective left $A$-module.

    Consider the natural map
    \[
        Af \longrightarrow \overline{A}e
    \]
    defined by
    \[
        af \longmapsto (a+J(A))(f+J(A)).
    \]
    Since $f+J(A)=e$, this map is explicitly
    \[
        af \longmapsto (a+J(A))e.
    \]
    It is a surjective $A$-module homomorphism. Its kernel is
    \[
        J(A)f.
    \]
    Indeed, $af$ maps to zero if and only if
    \[
        (a+J(A))(f+J(A))=0,
    \]
    which is equivalent to
    \[
        af\in J(A).
    \]
    Since $af\in Af$, the kernel is
    \[
        Af\cap J(A)=J(A)f.
    \]
    Thus
    \[
        Af/J(A)f\cong \overline{A}e.
    \]
    Because $\overline{A}e\cong S$, we get
    \[
        P_S/J(A)P_S\cong S.
    \]
    Since $A$ is Artinian and $P_S$ is finitely generated, we have
    \[
        \operatorname{Rad}(P_S)=J(A)P_S.
    \]
    Therefore
    \[
        P_S/\operatorname{Rad}(P_S)\cong S.
    \]

    Next we verify the stated property of $f$. Since every simple $A$-module is annihilated
    by $J(A)$, the action of $f$ on any simple $A$-module is the same as the action of
    \[
        e=f+J(A)
    \]
    through the quotient algebra $\overline{A}=A/J(A)$. Hence
    \[
        fS=eS\neq 0.
    \]
    If $T$ is a simple $A$-module with $T\not\cong S$, then, by the choice of $e$ in the
    semisimple algebra $\overline{A}$, we have
    \[
        eT=0.
    \]
    Therefore
    \[
        fT=0.
    \]

    Finally, the natural epimorphism
    \[
        P_S=Af \longrightarrow P_S/\operatorname{Rad}(P_S)\cong S
    \]
    is essential by Nakayama's lemma. Since $P_S$ is projective, this epimorphism is a
    projective cover of $S$.

    Thus $P_S$ is the projective cover of $S$.
\end{proof}

\begin{theorem}\label{thm:indecomposable-projectives-artinian}
    Let $A$ be an Artinian ring. Up to isomorphism, the indecomposable projective
    left $A$-modules are exactly the modules
    \[
        P_S,
    \]
    where $S$ runs through the isomorphism classes of simple left $A$-modules and
    $P_S$ denotes the projective cover of $S$.

    Moreover,
    \[
        P_S\cong P_T
        \qquad\Longleftrightarrow\qquad
        S\cong T.
    \]

    Each $P_S$ occurs as a direct summand of the left regular module ${}_A A$, with
    multiplicity equal to the multiplicity of $S$ as a direct summand of
    \[
        A/J(A).
    \]
    Equivalently, one has a decomposition
    \[
        {}_A A\cong \bigoplus_{S} P_S^{\,n_S},
    \]
    where $S$ runs through the isomorphism classes of simple left $A$-modules and
    \[
        A/J(A)\cong \bigoplus_S S^{\,n_S}.
    \]
    Furthermore,
    \[
        n_S=\dim_{D_S} S,
        \qquad
        D_S:=\operatorname{End}_A(S),
    \]
    where $S$ is regarded as a right $D_S$-vector space by
    \[
        s\cdot \varphi:=\varphi(s).
    \]
    In particular, if $A$ is a finite-dimensional $k$-algebra over an algebraically closed
    field $k$, then
    \[
        n_S=\dim_k S.
    \]
\end{theorem}

\begin{proof}
    We first classify the indecomposable projective modules.

    Let $P$ be a finitely generated indecomposable projective left $A$-module. Since $A$
    is Artinian, $P$ has finite length, and therefore
    \[
        P/\operatorname{Rad}(P)
    \]
    is semisimple. Hence we may write
    \[
        P/\operatorname{Rad}(P)\cong S_1\oplus \cdots \oplus S_t,
    \]
    where each $S_i$ is simple.

    The canonical epimorphism
    \[
        P\longrightarrow P/\operatorname{Rad}(P)
    \]
    is essential by Nakayama's lemma. Since $P$ is projective, this epimorphism is a
    projective cover of
    \[
        S_1\oplus \cdots \oplus S_t.
    \]
    For each simple left $A$-module $T$, let
    \[
        \pi_T:P_T\longrightarrow T
    \]
    denote the projective cover of $T$, whose existence was proved in
    Theorem~\ref{thm:projective-cover-simple}. Thus $P_T$ is an indecomposable
    projective left $A$-module and
    \[
        P_T/\operatorname{Rad}(P_T)\cong T.
    \]

    In particular, for each $i=1,\dots,t$, we have a projective cover
    \[
        \pi_{S_i}:P_{S_i}\longrightarrow S_i.
    \]
    Therefore the direct sum map
    \[
        \pi_{S_1}\oplus\cdots\oplus \pi_{S_t}:
        P_{S_1}\oplus \cdots \oplus P_{S_t}
        \longrightarrow
        S_1\oplus \cdots \oplus S_t
    \]
    is also a projective cover. By uniqueness of projective covers, we obtain
    \[
        P\cong P_{S_1}\oplus \cdots \oplus P_{S_t}.
    \]
    Since $P$ is indecomposable, we must have $t=1$. Hence
    \[
        P\cong P_{S_1}.
    \]
    Thus every indecomposable projective module is isomorphic to $P_S$ for some simple
    module $S$.

    Conversely, for every simple module $S$, the projective cover $P_S$ is indecomposable.
    Indeed, by the previous theorem, $P_S$ has the form
    \[
        P_S\cong Af
    \]
    for some primitive idempotent $f\in A$. Since $f$ is primitive, the module $Af$ is
    indecomposable.

    Next we show that
    \[
        P_S\cong P_T
        \qquad\Longleftrightarrow\qquad
        S\cong T.
    \]
    If $P_S\cong P_T$, then passing to radical quotients gives
    \[
        P_S/\operatorname{Rad}(P_S)
        \cong
        P_T/\operatorname{Rad}(P_T).
    \]
    But
    \[
        P_S/\operatorname{Rad}(P_S)\cong S,
        \qquad
        P_T/\operatorname{Rad}(P_T)\cong T,
    \]
    so $S\cong T$. The converse is immediate from the uniqueness of projective covers.

    Now consider the left regular module ${}_A A$. Since $A$ is Artinian, the natural
    epimorphism
    \[
        A\longrightarrow A/J(A)
    \]
    is essential by Nakayama's lemma. Since ${}_A A$ is projective, this epimorphism is
    a projective cover of $A/J(A)$.

    Write the semisimple $A$-module $A/J(A)$ as
    \[
        A/J(A)\cong \bigoplus_S S^{\,n_S},
    \]
    where $S$ runs through the isomorphism classes of simple left $A$-modules. Then
    \[
        \bigoplus_S P_S^{\,n_S}
        \longrightarrow
        \bigoplus_S S^{\,n_S}
    \]
    is a projective cover. By uniqueness of projective covers, we get
    \[
        {}_A A\cong \bigoplus_S P_S^{\,n_S}.
    \]
    Thus the multiplicity of $P_S$ in ${}_A A$ is exactly the multiplicity of $S$ in
    $A/J(A)$.

    It remains to identify $n_S$. Since $A/J(A)$ is semisimple Artinian, the Wedderburn--Artin
    decomposition shows that the multiplicity of a simple module $S$ in the left regular
    module of $A/J(A)$ is
    \[
        n_S=\dim_{D_S}S,
        \qquad
        D_S=\operatorname{End}_A(S).
    \]
    Here $S$ is regarded as a right $D_S$-vector space via
    \[
        s\cdot \varphi=\varphi(s).
    \]

    Finally, if $A$ is a finite-dimensional $k$-algebra over an algebraically closed field
    $k$, then by Schur's lemma,
    \[
        \operatorname{End}_A(S)\cong k.
    \]
    Therefore
    \[
        n_S=\dim_{D_S}S=\dim_k S.
    \]
\end{proof}

\begin{corollary}\label{cor:fg-module-has-projective-cover}
    Let $A$ be an Artinian ring, and let $U$ be a finitely generated left $A$-module.
    Then $U$ has a projective cover.
\end{corollary}

\begin{proof}
    Since $A$ is Artinian and $U$ is finitely generated, the quotient
    \[
        U/\operatorname{Rad}(U)
    \]
    is semisimple. Hence we may write
    \[
        U/\operatorname{Rad}(U)\cong S_1\oplus\cdots\oplus S_n,
    \]
    where each $S_i$ is a simple left $A$-module.

    For each $i$, let
    \[
        \pi_i:P_{S_i}\longrightarrow S_i
    \]
    be the projective cover of $S_i$. Then the direct sum
    \[
        \pi:=\pi_1\oplus\cdots\oplus\pi_n:
        P_{S_1}\oplus\cdots\oplus P_{S_n}
        \longrightarrow
        S_1\oplus\cdots\oplus S_n
    \]
    is a projective cover.

    Let
    \[
        q:U\longrightarrow U/\operatorname{Rad}(U)
    \]
    be the canonical epimorphism. After identifying
    \[
        U/\operatorname{Rad}(U)\cong S_1\oplus\cdots\oplus S_n,
    \]
    we may regard $\pi$ as an epimorphism
    \[
        \pi:P_{S_1}\oplus\cdots\oplus P_{S_n}
        \longrightarrow
        U/\operatorname{Rad}(U).
    \]

    Since
    \[
        P:=P_{S_1}\oplus\cdots\oplus P_{S_n}
    \]
    is projective and $q$ is surjective, there exists an $A$-module homomorphism
    \[
        \widetilde{\pi}:P\longrightarrow U
    \]
    such that the following diagram commutes:
    \[
        \begin{tikzcd}
            & P \arrow[d, "\pi"] \arrow[dl, dashed, "\widetilde{\pi}"'] \\
            U \arrow[r, "q"'] & U/\operatorname{Rad}(U) \arrow[r] & 0 .
        \end{tikzcd}
    \]
    That is,
    \[
        q\circ \widetilde{\pi}=\pi.
    \]

    We claim that $\widetilde{\pi}$ is an essential epimorphism. Since $\pi$ is surjective,
    $q\circ\widetilde{\pi}$ is surjective. Moreover, $q$ is essential by Nakayama's lemma.
    Therefore $\widetilde{\pi}$ is surjective.

    Now let $W\subseteq P$ be a submodule such that
    \[
        \widetilde{\pi}(W)=U.
    \]
    Then
    \[
        \pi(W)=q\bigl(\widetilde{\pi}(W)\bigr)=q(U)=U/\operatorname{Rad}(U).
    \]
    Since $\pi:P\to U/\operatorname{Rad}(U)$ is essential, it follows that
    \[
        W=P.
    \]
    Thus $\widetilde{\pi}$ is essential.

    Therefore
    \[
        \widetilde{\pi}:P_{S_1}\oplus\cdots\oplus P_{S_n}\longrightarrow U
    \]
    is an essential epimorphism from a projective module onto $U$. Hence it is a projective
    cover of $U$.
\end{proof}

\begin{lemma}\label{lem:artinian-noetherian}
    Let $A$ be a ring with identity.

    \begin{enumerate}
        \item If $A$ is semisimple, then $A$ is left Artinian if and only if $A$ is left
              Noetherian.

        \item If $A$ is left Artinian, then $A$ is left Noetherian.
    \end{enumerate}
\end{lemma}

\begin{proof}
    \textbf{(1)}
    Assume that $A$ is semisimple. Then the left regular module ${}_A A$ is semisimple,
    so it is a direct sum of simple left $A$-modules:
    \[
        {}_A A=\bigoplus_{\lambda\in \Lambda} S_\lambda .
    \]

    For a semisimple module, being Noetherian is equivalent to being a finite direct sum
    of simple modules, and the same is true for being Artinian. Indeed, if the index set
    $\Lambda$ is infinite, then we can form an infinite strictly increasing chain
    \[
        S_{\lambda_1}
        \subset
        S_{\lambda_1}\oplus S_{\lambda_2}
        \subset
        S_{\lambda_1}\oplus S_{\lambda_2}\oplus S_{\lambda_3}
        \subset \cdots ,
    \]
    so the module is not Noetherian. Similarly, we can form an infinite strictly decreasing
    chain by removing summands one by one, so the module is not Artinian.

    Thus a semisimple module is Noetherian if and only if it is Artinian, namely if and
    only if it is a finite direct sum of simple modules. Applying this to the left regular
    module ${}_A A$, we obtain that $A$ is left Artinian if and only if $A$ is left
    Noetherian.

    \medskip

    \textbf{(2)}
    Assume that $A$ is left Artinian. Let
    \[
        J=J(A)
    \]
    be the Jacobson radical of $A$. Since $A$ is left Artinian, the radical $J$ is nilpotent.
    Hence there exists $n\geq 1$ such that
    \[
        J^n=0.
    \]
    Therefore we have a finite filtration of left $A$-modules
    \[
        A\supseteq J\supseteq J^2\supseteq \cdots \supseteq J^n=0.
    \]

    It is enough to show that each quotient
    \[
        J^i/J^{i+1}
    \]
    is Noetherian as a left $A$-module. Since $J$ annihilates $J^i/J^{i+1}$, each quotient
    is naturally a module over
    \[
        A/J.
    \]
    Moreover, $A/J$ is semisimple Artinian. Hence every $A/J$-module is semisimple.

    Because $J^i/J^{i+1}$ is a quotient of the left Artinian module $J^i$, it is Artinian.
    Thus $J^i/J^{i+1}$ is an Artinian semisimple module. By part \textup{(1)}, it is
    Noetherian.

    Now using the finite filtration
    \[
        A\supseteq J\supseteq J^2\supseteq \cdots \supseteq J^n=0
    \]
    and the fact that a module is Noetherian whenever it has a finite filtration whose
    successive quotients are Noetherian, we conclude that ${}_A A$ is Noetherian.

    Therefore $A$ is left Noetherian.
\end{proof}

